% Propietario conservado: /Users/ruben/Documents/ChatGPT/jueces y controles/REVISION_ARGUMENTAL_APERTURAS_CIERRES_20260915/delivery/ARTICLE_V_ES_ARGUMENTAL_20260915/manuscrito/sections/v_coordenadas_angulares.tex
% RESULTADO_RECUPERADO desde II REV09. Véase MAPA_PROCEDENCIA.json.
% Selección de líneas y cambios mecánicos explícitos; ninguna prueba nueva.
\subsection{Coordenadas angulares, normalización de la vuelta y unidades de acción}
\label{v:v:ant:sec:ii-angulos}

La magnitud de acción y la coordenada angular pertenecen a espacios
distintos. La sección \(\hbar_{\rm ret}\) ocupa la recta de acción
\(\mathcal L_S\); \(\eta_{\rm ret}\) es su coeficiente en la base
declarada. La publicación angular actúa sobre la pareja adimensional
\((\alpha,\eta_{\rm ret})\), después de su construcción:
\begin{equation}
 \mathcal P(\alpha,\eta_{\rm ret})
 =\bigl([1000\alpha]_{360},[2(\eta_{\rm ret}+\alpha)]_{360}\bigr).
 \label{v:v:ant:eq:ii-par-angular}
\end{equation}
La completación nonádica determina el factor 1000; la clausura
bidireccional determina el factor dos. La carta circular de 360
unidades hace visibles las particiones
\(12\cdot30=4\cdot90=3\cdot120\). Sobre la rama positiva del
artículo, los representantes son
\[
 A_{\deg}=1000\alpha,\qquad
 C^*_{\deg}=2(\eta_{\rm ret}+\alpha).
\]

\begin{proposition}[Desplazamiento de la publicación en base mil]
\label{v:v:ant:prop:ii-desplazamiento-angular}
Sea \(\alpha=\sum_{n\geq1}b_n1000^{-n}\), con
\(b_n\in\{0,\ldots,999\}\), la expansión compatible ya generada.
La coordenada levantada \(A_{\deg}=1000\alpha\) satisface
\[
 A_{\deg}=b_1+\sum_{n\geq1}b_{n+1}1000^{-n}.
\]
La publicación conserva la sucesión de bloques a partir del segundo;
el bloque inicial pasa a la parte entera.
\end{proposition}
\begin{proof}
La serie converge absolutamente. Multiplicar cada término por 1000
y separar el término inicial proporciona la identidad. La operación
actúa sobre la expansión construida, con sus bloques y su memoria;
la publicación circular posterior toma su clase módulo 360.
\end{proof}

\begin{definition}[Cartas de la coordenada angular]
La carta en grados expresa una vuelta completa mediante 360 unidades.
La carta en radianes la expresa mediante \(2\pi\). El cambio de
coordenadas y sus inversas son
\[
 x=\frac{\pi}{180}A_{\deg},\qquad
 y=\frac{\pi}{180}C^*_{\deg},\qquad
 A_{\deg}=\frac{180}{\pi}x,\qquad
 C^*_{\deg}=\frac{180}{\pi}y.
\]
Sobre un recorrido con número de vueltas conservado se usan las
coordenadas levantadas; el cociente circular publica su clase módulo
la vuelta correspondiente.
\end{definition}

La definición caracteriza grados y radianes por su relación exacta
con la vuelta. El registro retiene además orientación, vueltas y
memoria de transporte. La publicación de su clase angular conserva
la fase, mientras el levantamiento distingue recorridos con la misma
fase final. La lectura física de acción necesita precisar cuál de
estas coordenadas interviene en la integral.

\begin{proposition}[Covariancia del lector de acción]
Para \(h_a=2\pi\hbar_a\), la acción de una coordenada levantada
satisface
\begin{equation}
 \mathcal S_a(\theta)
 =\hbar_a\theta_{\rm rad}
 =h_a\frac{\theta_{\deg}}{360}.
 \label{v:v:ant:eq:ii-accion-cartas}
\end{equation}
Una vuelta positiva publica \(h_a\).
\end{proposition}
\begin{proof}
Sustituir \(\theta_{\rm rad}=\pi\theta_{\deg}/180\) da
\(\hbar_a\pi\theta_{\deg}/180
=(2\pi\hbar_a)\theta_{\deg}/360\). En una vuelta
\(\theta_{\deg}=360\), de donde resulta \(h_a\).
\end{proof}

La periodicidad de una representación se declara aparte. Si la
representación transportada satisface \(U(2\pi)=-I\), su restitución
operatoria ocurre en \(4\pi\). Esta propiedad del levantamiento
coexiste con la normalización por vuelta de
\eqref{v:v:ant:eq:ii-accion-cartas}. Cada una conserva su objeto: magnitud
integrada, fase circular y estado representado.

\subsection{Covariancia de los canales y del funcional de incidencia}

Los canales se escriben de manera equivalente como
\begin{equation}
 q_\pm=\exp[-(x\pm y)]
 =\exp\!\left[-\frac{\pi}{180}(A_{\deg}\pm C^*_{\deg})\right].
\label{v:v:ant:eq:ii-canales}
\end{equation}
\begin{proposition}[Recuperación de las coordenadas par e impar]
\label{v:v:ant:prop:ii-inversion-canales-angulares}
En la carta real levantada, los canales positivos determinan
unívocamente las dos coordenadas:
\[
 A_{\deg}=-\frac{90}{\pi}\log(q_+q_-),\qquad
 C^*_{\deg}=\frac{90}{\pi}\log\frac{q_-}{q_+}.
\]
El intercambio \((q_+,q_-)\mapsto(q_-,q_+)\) conserva
\(A_{\deg}\) y revierte \(C^*_{\deg}\).
\end{proposition}
\begin{proof}
La suma y la diferencia de los logaritmos de
\eqref{v:v:ant:eq:ii-canales} son, respectivamente,
\(-\pi A_{\deg}/90\) y \(-\pi C^*_{\deg}/90\).
El logaritmo real es unívoco sobre los canales positivos.
El producto es simétrico y el cociente se invierte al intercambiarlos.
\end{proof}
En el registro de doce posiciones, la involución antipodal
\(r\mapsto r+6\pmod {12}\) forma seis parejas. La distribución
uniforme de la coordenada impar sobre estas parejas es
\(C^*_{\deg}/6\): sumar sus seis contribuciones recupera
\(C^*_{\deg}\). Esta coordenada sectorial se distingue de
la constante de Planck \(h\), que pertenece a la recta de acción.

La conversión se aplica una sola vez a cada argumento. La magnitud
\(\alpha\) sigue siendo adimensional y la magnitud \(\hbar\)
sigue perteneciendo a la recta de acción; sus coordenadas previas
intervienen en el mapa \(\mathcal P\).

El funcional de Barbero--Immirzi contiene, además de los canales, una
componente bilineal en los representantes angulares. Su transporte es
\begin{equation}
 \frac{\pi}{180}A_{\deg}C^*_{\deg}
 =\frac{180}{\pi}xy.
 \label{v:v:ant:eq:ii-barbero-cartas}
\end{equation}
La identidad se obtiene sustituyendo ambos representantes en grados
por sus expresiones en radianes. Por tanto, la fórmula completa en
coordenadas analíticas es
\[
 \gamma_{\HMT}
 =\frac{180}{\pi}xy
 +12\bigl[S_{90}(e^{-x+y})-S_{90}(e^{-x-y})\bigr]
 -\bigl[S_{120}(e^{-x+y})-S_{120}(e^{-x-y})\bigr].
\]
La equivalencia de las dos cartas transporta la componente bilineal
y las funciones exponenciales simultáneamente. Esta transformación
explicita el papel preciso del grado en la construcción y permite
reproducir la evaluación sin ambigüedad de unidades angulares.
